% !PRINTED 223
\weberpage{223}

\begin{center}
\S.\,72.
\end{center}

\begin{center}
\gesperrt{Allgemeine Ausführung der Transformation.}
\end{center}

In der allgemeinen Durchführung der \gesperrt{Tschirnhausen}-Transformation in der \gesperrt{Hermite'schen} Form können wir noch einen bedeutenden Schritt weiter gehen.

Wir betrachten zunächst die allgemeine Substitution \S.\,68, (5)
\weq{1}{
y = t_{n-1} f_0 + t_{n-2} f_1 + t_{n-3} f_2 + \cdots + t_0 f_{n-1},
}
aus der wir die speciellere Form \S.\,68 (12) erhalten, wenn wir die Variablen $t$ an die eine lineare Bedingung
\weq{2}{
S(y) = n a_0 t_{n-1} + (n-1) a_1 t_{n-2} + \cdots + a_{n-1} t_0 = 0
}
binden, was aber fürs Erste noch nicht geschehen soll.

Durch die Functionen $f_0, f_1, f_2 \ldots f_{n-1}$ lassen sich nach \S.\,68, (4) alle Potenzen von $x$ und also auch alle rationalen Functionen von $x$ \gesperrt{linear} und homogen darstellen.

Wenn wir diese Darstellung für die verschiedenen Potenzen von $y$ finden können, so lässt sich durch Elimination der $f_i$ die Gleichung $n^{\text{ten}}$ Grades für $y$ bilden.

Denselben Zweck erreichen wir aber noch einfacher, wenn wir die $n$ Producte $y f_s$ linear durch die $f$ darstellen. Denn wenn wir die Gleichungen haben
\weq{3}{
y f_s = E_{0,s} f_0 + E_{1,s} f_1 + E_{2,s} f_2 + \cdots + E_{n-1,s} f_{n-1}
}
\[
s = 0, 1, 2 \ldots n-1,
\]
so erhalten wir durch Elimination von $f_0, f_1 \ldots f_{n-1}$ aus dem System (3)
\weq{4}{
E =
\begin{vmatrix}
E_{0,0} - y, & E_{1,0}, & E_{2,0} & \cdots & E_{n-1,0} \\
E_{0,1}, & E_{1,1} - y, & E_{2,1} & \cdots & E_{n-1,1} \\
\cdots & \cdots & \cdots & \cdots & \cdots \\
E_{0,n-1}, & E_{1,n-1}, & E_{2,n-1} & \cdots & E_{n-1,n-1} - y
\end{vmatrix}
= 0,
}
also eine Gleichung $n^{\text{ten}}$ Grades für $y$, die die gesuchte transformirte Gleichung ist. Alles ist daher zurückgeführt auf die Bestimmung der Co\"efficienten $E_{r,s}$, die, wie aus (1) hervorgeht, jedenfalls lineare Functionen der Variablen $t$ sind. Insbesondere ist
\weq{5}{
E_{0,0} = a_0 t_{n-1}, \quad E_{1,0} = a_0 t_{n-2}, \quad \ldots \quad E_{n-1,0} = a_0 t_0.
}

Um die übrigen $E$ zu berechnen, bemerken wir, dass nach der Definition \S.\,68, (4) zwischen den Functionen $f$ die folgenden Relationen bestehen:
%CONT>>
